\clearpage 

\section {Summary}
\label{sec:summary}

This technical note reports a joint 3-flavour oscillation analysis by the VaLOR group on the combined Run 1+2+3 dataset, 
which corresponds to an integrated neutrino beam exposure of 3.020$\times$10$^{20}$ POT. 
The joint analysis determines $|\Delta m^{2}|$, $\sin^{2} \theta_{23}$,  $\sin^{2} \theta_{13}$ and $\delta_{CP}$ together 
by making a simultaneous fit of the reconstructed energy spectra of the single $\mu$-like ring and single e-like ring SuperK event samples. 
This fit is performed in a 3-flavour framework including constant-density matter effects. Official 2012 T2K analysis inputs are used, 
and all 83 systematic parameters considered in the analysis (see Sec.\ \ref{sec:systematics}) are allowed to float in the fit. 
The Fogli-Lisi convention is used for the ``atmospheric'' mass-squared splitting $|\Delta m^{2}|$ (see Sec.\ \ref{sec:oscparamchoice}). 
The ``solar'' oscillation parameters $\sin^{2} \theta_{12}$ and $\Delta m_{21}^{2}$ have been found to have a negligible effect (see appendix \ref{app:solar}), 
and are fixed at 0.311 and 7.5$\times$10$^{-5}$ eV$^{2}$/c$^{4}$ respectively (the 2012 PDG values). 

This analysis has been extensively validated. The $\nu_{\mu}$-disappearance analysis is identical to the 2012 VaLOR analysis reported in Ref.\ \cite{OurRun123DisappAnalysis3Flavour}, 
and the $\nu_{e}$-appearance analysis has been validated against the official 2012 analyses (see appendix \ref{app:standalone}). 
Studies have been made of fitter residuals (appendix \ref{app:fitter_performance}), fits of 3.020$\times$10$^{20}$ POT fake datasets (appendix \ref{app:fake_dataset_fits}) 
and the T2K expected sensitivity for 3.020$\times$10$^{20}$ POT (appendix \ref{app:sensitivity}). These studies have shown that the joint analysis fitter is performing as expected. 

In the present version of the analysis, correlations between the single $\mu$-like ring and single e-like ring systematic uncertainties have been neglected, 
but they will be studied and reported in a future version. 

For the normal hierarchy (NH), the joint analysis fit of the reconstructed energy spectra of the single $\mu$-like ring and single e-like ring SuperK event samples gives
$|\Delta m^{2}|$ = 2.494$\times$10$^{-3}$ $eV^{2}/c^{4}$,
$\sin^{2}\theta_{23}$ = 0.516,
$\sin^{2}\theta_{13}$ = 0.029 and
$\delta_{CP}$ = 2.086.
The corresponding best-fit values for the inverted hierarchy (IH) are 
$|\Delta m^{2}|$ = 2.468$\times$10$^{-3}$ $eV^{2}/c^{4}$,
$\sin^{2}\theta_{23}$ = 0.516,
$\sin^{2}\theta_{13}$ = 0.034 and
$\delta_{CP}$ = 0.922.

Confidence regions are constructed using the constant-$\Delta\chi^2$ method.
For the NH, the 90\% (68\%) C.L. allowed region for the four fitted oscillation parameters can be summarized as
2.227$\times$10$^{-3}$ (2.332$\times$10$^{-3}$) eV$^{2}$/c$^{4}$
$< |\Delta m^{2}| < $
2.773$\times$10$^{-3}$ (2.659$\times$10$^{-3}$) eV$^{2}$/c$^{4}$,
0.401 (0.433) $< \sin^{2}\theta_{23} <$ 0.628 (0.596),
0.006 (0.012) $< \sin^{2}\theta_{13} <$ 0.058 (0.048) and
-$\pi$ (-$\pi$) $< \delta_{CP} <$ $\pi$ ($\pi$).
For the IH, the corresponding 90\% (68\%) C.L. allowed regions are
2.218 $\times$10$^{-3}$ (2.311$\times$10$^{-3}$) eV$^{2}$/c$^{4}$
$< |\Delta m^{2}| < $
2.729 $\times$10$^{-3}$ (2.627$\times$10$^{-3}$) eV$^{2}$/c$^{4}$,
0.397 (0.433) $< \sin^{2}\theta_{23} <$ 0.627 (0.596),
0.006 (0.013) $< \sin^{2}\theta_{13} <$ 0.069 (0.054) and
-$\pi$ (-$\pi$) $< \delta_{CP} <$ $\pi$ ($\pi$).

The joint analysis has the advantage compared with the stand-alone analyses that it does not need to assume the values of any of the oscillation 
parameters $|\Delta m^{2}|$, $\sin^{2} \theta_{23}$,  $\sin^{2} \theta_{13}$ and $\delta_{CP}$. In particular, the value of $\sin^{2} \theta_{23}$ 
is not assumed as in the stand-alone $\nu_{e}$-appearance analysis, but is constrained by the fit 
of the reconstructed energy spectrum of the single $\mu$-like ring SuperK event sample.

